% !TEX program = xelatex
\documentclass[11pt,a4paper]{article}
\usepackage{algo-preamble}

\title{\textbf{L10 \textperiodcentered{} Δομές Δεδομένων}\\[2pt]
\large Σωροί \& Ξένα Σύνολα}
\author{Algorithms Class Hub}
\date{\small Αλγόριθμοι και Πολυπλοκότητα (K17) \textperiodcentered{} ΕΚΠΑ}

\begin{document}
\maketitle

\begin{center}\itshape\small
Δυαδικοί σωροί (heaps), heapify-up / heapify-down, ουρά προτεραιότητας,
heapsort, και η δομή ξένων συνόλων (union-find) με βάρη στα δέντρα και συμπίεση
μονοπατιών.
\end{center}

\section{Τι θα δούμε}

Στο προηγούμενο μάθημα είπαμε ότι ο \textbf{Dijkstra} και ο \textbf{Prim} τρέχουν
σε $O(m \log n)$ «με δυαδικό σωρό», και ο \textbf{Kruskal} χρειάζεται μια δομή
που απαντάει «είναι αυτές οι δύο κορυφές στην ίδια συνεκτική συνιστώσα;». Είπαμε
«θα το δούμε αργότερα». \textbf{Αυτή είναι η διάλεξη του «αργότερα».}

Δύο δομές δεδομένων:
\begin{enumerate}
  \item \textbf{Σωρός (heap)} --- υλοποιεί μια \textbf{ουρά προτεραιότητας}:
  «δώσε μου το μικρότερο στοιχείο» σε $O(\log n)$. Αυτό κάνει τον Dijkstra και
  τον Prim γρήγορους.
  \item \textbf{Ξένα σύνολα (union-find)} --- υλοποιεί «σε ποιο σύνολο ανήκει το
  $x$;» και «ένωσε δύο σύνολα» σχεδόν σε σταθερό χρόνο. Αυτό κάνει τον Kruskal
  γρήγορο.
\end{enumerate}

\begin{keypoint}
\textbf{Δεν είναι «θεωρία για τη θεωρία».} Κάθε δομή εδώ υπάρχει επειδή ένας
αλγόριθμος που ήδη ξέρεις την χρειάζεται. Διάβασε κάθε ενότητα ρωτώντας: «ποια
ακριβώς πράξη επιταχύνει αυτό, και σε ποιον αλγόριθμο;»
\end{keypoint}

\section{Μέρος Α' --- Σωροί (Heaps)}

\subsection{Το πρόβλημα που λύνουν}

Φαντάσου ότι έχεις μια συλλογή στοιχείων με \textbf{κλειδιά} (αριθμούς) και
θέλεις, ξανά και ξανά, να ρωτάς: \textbf{«ποιο είναι το μικρότερο;»} και μετά να
το \textbf{βγάζεις} από τη συλλογή. Ακριβώς αυτό κάνει ο Dijkstra: σε κάθε βήμα
βγάζει την ανεξερεύνητη κορυφή με το μικρότερο $\pi(v)$.

\begin{center}
\begin{tabular}{@{}l l l@{}}
\toprule
Δομή & FindMin & Insert / Delete \\
\midrule
Μη ταξινομημένος πίνακας & $O(n)$ --- σάρωσε όλα      & $O(1)$ \\
Ταξινομημένος πίνακας    & $O(1)$ --- είναι στην άκρη & $O(n)$ --- μετατόπιση \\
Σωρός                    & $O(1)$                     & $O(\log n)$ \\
\bottomrule
\end{tabular}
\end{center}

\begin{intuition}
\textbf{Διαισθητικά:} ο σωρός είναι ο \textbf{συμβιβασμός}. Ο ταξινομημένος
πίνακας κρατάει τα πάντα σε τέλεια σειρά --- υπερβολικά πολλή δουλειά. Ο μη
ταξινομημένος δεν κρατάει τίποτα --- πληρώνεις στο ψάξιμο. Ο σωρός κρατάει
\textbf{μόνο όση τάξη χρειάζεται} για να ξέρει πάντα ποιο είναι το μικρότερο:
κάθε γονιός μικρότερος από τα παιδιά του. Τίποτα παραπάνω.
\end{intuition}

\subsection{Ορισμός}

\begin{keypoint}
\textbf{Σωρός (heap):} ένα \textbf{ισοσταθμισμένο δυαδικό δέντρο} που ικανοποιεί
την \textbf{ιδιότητα διάταξης σωρού}: για κάθε κορυφή $v$ με γονέα $w$ ισχύει
$\text{key}(w) \le \text{key}(v)$.
\end{keypoint}

Δύο κομμάτια, ξεχωριστά:
\begin{itemize}
  \item \textbf{«Ισοσταθμισμένο δυαδικό δέντρο»} --- αφορά το \textbf{σχήμα}.
  Κάθε επίπεδο γεμίζει πλήρως πριν ξεκινήσει το επόμενο, και το τελευταίο
  γεμίζει από αριστερά. Αυτό εγγυάται ύψος $\lfloor \log_2 n \rfloor$ --- εκεί
  κρύβεται το $O(\log n)$ κάθε πράξης.
  \item \textbf{«Ιδιότητα διάταξης σωρού»} --- αφορά τα \textbf{κλειδιά}. Κάθε
  γονιός $\le$ κάθε παιδί του.
\end{itemize}

Πρόσεξε τι \textbf{δεν} λέει η ιδιότητα: δεν λέει τίποτα για τη σχέση ανάμεσα σε
\textbf{αδέλφια}, ούτε ανάμεσα σε κορυφές διαφορετικών κλάδων. Ένας σωρός είναι
«χαλαρά» ταξινομημένος --- και ακριβώς γι' αυτό συντηρείται φθηνά.

\begin{keypoint}
\textbf{Συνέπεια:} η ρίζα έχει πάντα το \textbf{ελάχιστο} κλειδί. (Είναι γονιός
--- άμεσα ή έμμεσα --- όλων.) Άρα \texttt{FindMin} $= O(1)$: κοίτα τη ρίζα.
\end{keypoint}

\subsection{Αναπαράσταση με πίνακα}

Το ωραιότερο κόλπο: ένας σωρός \textbf{δεν χρειάζεται δείκτες}. Επειδή το σχήμα
είναι αυστηρά ισοσταθμισμένο, αποθηκεύουμε τις κορυφές σε έναν πίνακα $H[1..n]$,
επίπεδο-επίπεδο από αριστερά προς δεξιά. Τότε για κάθε θέση $i$:
\[
\text{parent}(i) = \left\lfloor i/2 \right\rfloor, \qquad
\text{leftChild}(i) = 2i, \qquad
\text{rightChild}(i) = 2i + 1
\]
Δηλαδή «να κατέβεις στο παιδί» $=$ πολλαπλασίασε με 2· «να ανέβεις στον γονέα»
$=$ διαίρεσε με 2. Καθαρή αριθμητική, μηδέν δείκτες.

\textbf{Παράδειγμα.} Ένας σωρός με κλειδιά $H = [2, 5, 8, 9, 11, 14, 18]$: η ρίζα
$H[1] = 2$ έχει παιδιά $H[2] = 5$ και $H[3] = 8$· το $H[2] = 5$ έχει παιδιά
$H[4] = 9$ και $H[5] = 11$. Κάθε γονιός είναι μικρότερος από τα παιδιά του
($2 \le 5, 8$· $5 \le 9, 11$· $8 \le 14, 18$), αλλά τα αδέλφια δεν είναι
ταξινομημένα μεταξύ τους --- και δεν χρειάζεται.

\subsection{Heapify-up --- επιδιόρθωση προς τη ρίζα}

Όταν προσθέτουμε ένα στοιχείο, το βάζουμε στο \textbf{τέλος} του πίνακα (πρώτη
ελεύθερη θέση --- διατηρεί το σχήμα). Όμως το κλειδί του μπορεί να είναι
\textbf{πολύ μικρό} και να παραβιάζει την ιδιότητα: μπορεί να είναι μικρότερο από
τον γονέα του. Η \texttt{Heapify-up} το διορθώνει \textbf{σπρώχνοντάς το προς τα
πάνω}:

\begin{codebox}
\begin{verbatim}
Heapify-up(H, i):
  Εάν i > 1 τότε:
    j <- parent(i) = ⌊i/2⌋
    Εάν key(H[i]) < key(H[j]) τότε:
      Εναλλαγή των H[i] και H[j]
      Heapify-up(H, j)
\end{verbatim}
\end{codebox}

\begin{intuition}
\textbf{Διαισθητικά:} το νέο, πολύ ελαφρύ στοιχείο «αναδύεται» σαν φυσαλίδα προς
την επιφάνεια. Σε κάθε βήμα ανταλλάσσεται με τον γονέα του, μέχρι ή να βρει γονέα
μικρότερο από αυτό, ή να γίνει το ίδιο ρίζα.
\end{intuition}

\textbf{Γιατί διορθώνει σωστά:} όταν ανταλλάσσουμε το $H[i]$ με τον γονέα του, το
μικρό κλειδί ανεβαίνει --- και ο γονέας, που ήταν μεγαλύτερος, κατεβαίνει σε μια
θέση όπου τα παιδιά του ήταν ήδη $\ge$ από το παλιό περιεχόμενό της, άρα $\ge$
από αυτόν. Δεν δημιουργείται νέα παραβίαση πουθενά.

\textbf{Χρόνος:} το μονοπάτι προς τη ρίζα έχει μήκος $\le \log_2 n$, και κάθε
βήμα είναι $O(1)$. Άρα \texttt{Heapify-up} $= O(\log n)$.

\textbf{Χρήση --- Insert:} τοποθέτησε το νέο στοιχείο στο τέλος, κάλεσε
\texttt{Heapify-up}. Συνολικά $O(\log n)$.

\subsection{Heapify-down --- επιδιόρθωση προς τα φύλλα}

Συμμετρικό πρόβλημα: μια κορυφή έχει κλειδί \textbf{πολύ μεγάλο}, μεγαλύτερο από
κάποιο παιδί της. Τη σπρώχνουμε \textbf{προς τα κάτω}:

\begin{codebox}
\begin{verbatim}
Heapify-down(H, i):
  Έστω l = 2i, r = 2i + 1   (παιδιά της i)
  Έστω j η θέση του παιδιού με το ΜΙΚΡΟΤΕΡΟ κλειδί (αν υπάρχουν παιδιά)
  Εάν υπάρχει j και key(H[j]) < key(H[i]) τότε:
    Εναλλαγή των H[i] και H[j]
    Heapify-down(H, j)
\end{verbatim}
\end{codebox}

\begin{warningbox}
\textbf{Κρίσιμη λεπτομέρεια:} ανταλλάσσουμε με το \textbf{μικρότερο} από τα δύο
παιδιά --- όχι με όποιο τύχει. Αν ανταλλάσσαμε με το μεγαλύτερο, αυτό θα
κατέληγε γονιός του μικρότερου αδελφού του και θα παραβιαζόταν ξανά η ιδιότητα.
Με το μικρότερο, ο νέος γονιός είναι σίγουρα $\le$ και από τα δύο παιδιά.
\end{warningbox}

\textbf{Χρόνος:} ίδιο επιχείρημα --- μονοπάτι προς τα φύλλα μήκους
$\le \log_2 n$. \texttt{Heapify-down} $= O(\log n)$.

\subsection{Ουρά προτεραιότητας}

Συνδυάζοντας τα παραπάνω παίρνουμε την \textbf{ουρά προτεραιότητας} --- μια δομή
που πάντα ξέρει «τι έχει σειρά»:

\begin{center}
\begin{tabular}{@{}l p{6.2cm} l@{}}
\toprule
Πράξη & Τι κάνει & Χρόνος \\
\midrule
\texttt{CreateHeap(H)} & κενός σωρός για $\le n$ στοιχεία & $O(n)$ \\
\texttt{Insert(H, v)}  & βάλε στο τέλος, \texttt{Heapify-up} & $O(\log n)$ \\
\texttt{FindMin(H)}    & επίστρεψε τη ρίζα $H[1]$ & $O(1)$ \\
\texttt{Delete(H, i)}  & βγάλε τη θέση $i$, \texttt{Heapify-down} (ή up) &
$O(\log n)$ \\
\texttt{ExtractMin(H)} & \texttt{FindMin} $+$ \texttt{Delete} της ρίζας &
$O(\log n)$ \\
\bottomrule
\end{tabular}
\end{center}

\begin{intuition}
\textbf{Πώς δουλεύει το \texttt{ExtractMin}:} η απάντηση είναι η ρίζα ---
εύκολο. Αλλά τώρα η ρίζα είναι «τρύπα». Φέρνουμε το \textbf{τελευταίο} στοιχείο
του πίνακα στη ρίζα (διατηρεί το σχήμα), και κάνουμε \texttt{Heapify-down} ---
αυτό το (μάλλον μεγάλο) στοιχείο βυθίζεται στη σωστή του θέση.
\end{intuition}

Η ουρά προτεραιότητας υποστηρίζει επίσης \texttt{decreaseKey(H, v, k)}: μειώνει
το κλειδί ενός στοιχείου και καλεί \texttt{Heapify-up}. Αυτή ακριβώς είναι η
πράξη που κάνουν ο Dijkstra και ο Prim όταν βρίσκουν φθηνότερο μονοπάτι/ακμή προς
μια κορυφή.

\subsection{Εφαρμογή --- Heapsort}

Με την ουρά προτεραιότητας παίρνουμε δωρεάν έναν αλγόριθμο ταξινόμησης:
\begin{enumerate}
  \item \texttt{CreateHeap} και βάλε μέσα και τα $n$ στοιχεία.
  \item Κάλεσε \texttt{ExtractMin} $n$ φορές --- βγαίνουν σε \textbf{αύξουσα}
  σειρά.
\end{enumerate}
\textbf{Χρόνος:} $n$ φορές το $O(\log n)$ $\to$ \textbf{$O(n \log n)$}. Είναι το
ίδιο φράγμα με τη συγχωνευτική ταξινόμηση του τρίτου μαθήματος, αλλά με εντελώς
διαφορετική ιδέα.

\subsection{Πώς οι σωροί επιταχύνουν τον Dijkstra και τον Prim}

Τώρα κλείνει η υπόσχεση του προηγούμενου μαθήματος. Και οι δύο αλγόριθμοι
κρατάνε μια ουρά προτεραιότητας από ανεξερεύνητες κορυφές:
\begin{itemize}
  \item \texttt{ExtractMin} --- διάλεξε την επόμενη κορυφή προς εξερεύνηση.
  Καλείται $n$ φορές $\to O(n \log n)$.
  \item \texttt{decreaseKey} --- όταν μια κορυφή $v$ μπει στο $S$, για κάθε ακμή
  $(v, w)$ ενημέρωσε την προτεραιότητα του $w$. Συνολικά μέχρι $m$ φορές $\to
  O(m \log n)$.
\end{itemize}
Σύνολο: \textbf{$O(m \log n)$}. Χωρίς σωρό, το \texttt{ExtractMin} θα ήταν
γραμμική σάρωση $O(n)$ και θα κατέληγες σε $O(n^2)$.

\section{Μέρος Β' --- Ξένα Σύνολα (Union-Find)}

\subsection{Το πρόβλημα που λύνουν}

Θυμήσου τον \textbf{Kruskal} από το προηγούμενο μάθημα: σαρώνει ακμές σε αύξουσα
σειρά κόστους και προσθέτει την $\{u, v\}$ \textbf{εκτός αν δημιουργεί κύκλο}.
Πώς ελέγχεις «δημιουργεί κύκλο;» Μια ακμή $\{u, v\}$ κλείνει κύκλο
\textbf{ακριβώς όταν} οι $u$ και $v$ είναι \textbf{ήδη} στην ίδια συνεκτική
συνιστώσα του μέχρι τώρα δέντρου.

Άρα χρειαζόμαστε μια δομή που κρατάει μια \textbf{διαμέριση} --- μια συλλογή από
\textbf{ανά δύο ξένα σύνολα} --- και απαντά γρήγορα: «το $u$ και το $v$ είναι στο
ίδιο σύνολο;»

\subsection{Οι τρεις πράξεις}

\begin{keypoint}
Η δομή \textbf{ξένων συνόλων (union-find)} διατηρεί μια συλλογή
$\mathcal{S} = \{S_1, \dots, S_k\}$ από ανά δύο ξένα σύνολα. Κάθε σύνολο έχει
έναν \textbf{αντιπρόσωπο} (ένα μέλος του). Πράξεις:
\begin{itemize}
  \item \texttt{MakeSet(x)} --- φτιάξε νέο σύνολο με μοναδικό μέλος (και
  αντιπρόσωπο) το $x$.
  \item \texttt{Union(x, y)} --- συγχώνευσε τα σύνολα που περιέχουν τα $x$ και
  $y$ σε ένα.
  \item \texttt{Find(x)} --- επίστρεψε τον αντιπρόσωπο του συνόλου που περιέχει
  το $x$.
\end{itemize}
\end{keypoint}

Ο αντιπρόσωπος είναι το «όνομα» του συνόλου: δύο στοιχεία είναι στο ίδιο σύνολο
\textbf{αν και μόνο αν} έχουν τον ίδιο \texttt{Find}.

\subsection{Εφαρμογή --- συνεκτικές συνιστώσες}

Με αυτές τις τρεις πράξεις, η εύρεση των συνεκτικών συνιστωσών ενός γραφήματος
γίνεται μηχανικά:

\begin{codebox}
\begin{verbatim}
ΣυνεκτικέςΣυνιστώσες(G = (V, E)):
  Για κάθε κορυφή v ∈ V:
    MakeSet(v)
  Για κάθε ακμή {u, v} ∈ E:
    Εάν Find(u) ≠ Find(v):
      Union(u, v)
\end{verbatim}
\end{codebox}

Στην αρχή κάθε κορυφή είναι μόνη της. Κάθε ακμή «κολλάει» τα δύο της άκρα στην
ίδια συνιστώσα. Στο τέλος, για δύο κορυφές $x, y$, ισχύει «ίδια συνιστώσα» αν και
μόνο αν $\texttt{Find}(x) = \texttt{Find}(y)$.

\subsection{Μια αφελής προσέγγιση}

Η πιο απλή ιδέα: ένας πίνακας $\text{label}[\,]$ όπου $\text{label}[x] = $ το
«όνομα» του συνόλου του $x$.
\begin{itemize}
  \item \texttt{Find(x)} --- επίστρεψε $\text{label}[x]$. \textbf{$O(1)$} ---
  εξαιρετικό.
  \item \texttt{Union(x, y)} --- διάλεξε ένα όνομα και \textbf{σάρωσε όλον τον
  πίνακα} αλλάζοντας κάθε εμφάνιση του άλλου ονόματος. \textbf{$O(n)$} ---
  απαίσιο.
\end{itemize}
Με $n$ ενώσεις φτάνουμε $O(n^2)$. Πληρώνουμε ακριβά στο \texttt{Union}.
Χρειαζόμαστε καλύτερη ισορροπία.

\subsection{Κατευθυνόμενα δέντρα με ρίζα}

Η σωστή ιδέα: αναπαριστούμε \textbf{κάθε σύνολο σαν ένα κατευθυνόμενο δέντρο με
ρίζα}.
\begin{itemize}
  \item Κάθε κορυφή έχει \textbf{έναν} δείκτη: προς τον \textbf{γονέα} της.
  \item Η \textbf{ρίζα} δείχνει στον εαυτό της (θηλιά) --- και είναι ο
  \textbf{αντιπρόσωπος} του συνόλου.
  \item Πολλά ξένα σύνολα $\to$ ένα κατευθυνόμενο \textbf{δάσος} με ρίζες.
\end{itemize}

\subsection{Υλοποίηση των πράξεων}

\begin{itemize}
  \item \texttt{MakeSet(x)} --- φτιάξε δέντρο με μία κορυφή $x$ (θηλιά στον
  εαυτό της). $O(1)$.
  \item \texttt{Find(x)} --- ακολούθησε τους δείκτες-γονείς από το $x$ μέχρι να
  φτάσεις στη ρίζα (την κορυφή με θηλιά). Επίστρεψέ την.
  \item \texttt{Union(x, y)} --- βρες τις δύο ρίζες· κάνε τη ρίζα του ενός
  δέντρου να δείχνει στη ρίζα του άλλου (αφαιρώντας τη θηλιά της). $O(1)$
  \textbf{αφού} έχεις τις ρίζες.
\end{itemize}

\begin{warningbox}
\textbf{Το πρόβλημα:} το κόστος είναι ολόκληρο στο \texttt{Find}. Αν ένα δέντρο
εκφυλιστεί σε \textbf{μονοπάτι} (αλυσίδα), το \texttt{Find} από το φύλλο
διασχίζει $O(n)$ κορυφές. Πίσω στο τετράγωνο. Χρειαζόμαστε τα δέντρα να μένουν
\textbf{κοντά}.
\end{warningbox}

\subsection{Βάρη στα δέντρα (union by size)}

Πρώτη βελτίωση: για κάθε δέντρο αποθηκεύουμε το \textbf{μέγεθός} του, και στο
\texttt{Union(x, y)} κάνουμε \textbf{τη ρίζα του μικρότερου δέντρου να δείχνει
στη ρίζα του μεγαλύτερου} (όχι ανάποδα).

\begin{keypoint}
\textbf{Θεώρημα.} Με union-by-size, σε ένα σύνολο με $k$ στοιχεία κάθε στοιχείο
απέχει το πολύ $\log_2 k$ βήματα από τη ρίζα του. Άρα \texttt{Find} $=
O(\log n)$, και αρκούν $2\log n$ βήματα για να αποφασίσουμε αν δύο στοιχεία είναι
στο ίδιο σύνολο.
\end{keypoint}

\textbf{Απόδειξη (επαγωγή στο πλήθος ενώσεων).} Ισχυρισμός: δέντρο με $k$ κορυφές
έχει βάθος $\le \log_2 k$.

\textit{Βάση:} δέντρο με $1$ κορυφή έχει βάθος $0 = \log_2 1$. $\checkmark$

\textit{Επαγωγικό βήμα:} ένα δέντρο με $k$ κορυφές προέκυψε από \texttt{Union}
δύο δέντρων με $k_1$ και $k_2$ κορυφές, $k_1 \le k_2$, $k_1 + k_2 = k$. Βάλαμε το
μικρό ($k_1$) κάτω από το μεγάλο ($k_2$). Το βάθος του νέου δέντρου είναι το
πολύ:
\[
\max\bigl\{\, \underbrace{\log_2 k_2}_{\text{το μεγάλο μένει ίδιο}},\ \
\underbrace{1 + \log_2 k_1}_{\text{το μικρό κατεβαίνει 1}} \,\bigr\}
\]
Αφού $k_1 \le k/2$, έχουμε $1 + \log_2 k_1 \le 1 + \log_2(k/2) = \log_2 k$. Και
$\log_2 k_2 \le \log_2 k$. Άρα το βάθος $\le \log_2 k$. $\blacksquare$

\begin{intuition}
\textbf{Γιατί δουλεύει:} μια κορυφή «κατεβαίνει» ένα επίπεδο μόνο όταν το δέντρο
της ενώνεται κάτω από ένα \textbf{τουλάχιστον ισομέγεθες}. Άρα κάθε φορά που
κατεβαίνει, το σύνολό της \textbf{τουλάχιστον διπλασιάζεται}. Δεν μπορεί να
διπλασιαστεί πάνω από $\log_2 n$ φορές $\to$ βάθος $\le \log_2 n$.
\end{intuition}

\subsection{Συμπίεση μονοπατιών (path compression)}

Δεύτερη βελτίωση, σχεδόν δωρεάν. Κάθε φορά που τρέχουμε \texttt{Find(v)} έχουμε
ήδη διασχίσει όλο το μονοπάτι από το $v$ μέχρι τη ρίζα $x$. Πριν επιστρέψουμε,
\textbf{ξανασυνδέουμε κάθε κορυφή αυτού του μονοπατιού απευθείας στη ρίζα $x$}.

\begin{intuition}
\textbf{Διαισθητικά:} «αφού πλήρωσα ήδη να ανέβω αυτό το μονοπάτι, ας το ισιώσω
για την επόμενη φορά.» Η επόμενη \texttt{Find} σε οποιαδήποτε από αυτές τις
κορυφές θα είναι ένα βήμα. Μικραίνουμε το \textbf{ύψος} των δέντρων χωρίς να
πειράζουμε τα μεγέθη/βάρη.
\end{intuition}

Με \textbf{union-by-size $+$ path compression} μαζί, μια ακολουθία $m$ πράξεων
σε $n$ στοιχεία τρέχει σε $O(m \cdot \alpha(n))$, όπου $\alpha$ είναι η
\textbf{αντίστροφη συνάρτηση Ackermann} --- μεγαλώνει τόσο αφάνταστα αργά που για
κάθε $n$ που θα συναντήσεις ποτέ στο σύμπαν, $\alpha(n) \le 4$. Πρακτικά
\textbf{σταθερά}.

\subsection{Πώς το union-find επιταχύνει τον Kruskal}

Κλείνει και η δεύτερη υπόσχεση του προηγούμενου μαθήματος. Ο Kruskal:

\begin{codebox}
\begin{verbatim}
Kruskal(G, c):
  Ταξινόμησε τις ακμές κατά κόστος          -> O(m log n)
  Για κάθε u ∈ V: MakeSet(u)
  Για i = 1 έως m, με {u, v} = e_i:
    Εάν Find(u) ≠ Find(v):                  <- «κλείνει κύκλο;»
      πρόσθεσε e_i στο T
      Union(u, v)
\end{verbatim}
\end{codebox}

\begin{itemize}
  \item \textbf{Ταξινόμηση:} $O(m \log n)$.
  \item \textbf{$m$ ζεύγη \texttt{Find} $+$ έως $n-1$ \texttt{Union}:}
  $O(m \cdot \alpha(n))$ --- ουσιαστικά γραμμικό.
  \item \textbf{Σύνολο:} $O(m \log n)$ --- η ταξινόμηση κυριαρχεί.
\end{itemize}

\begin{recapbox}
\textbf{Τι κρατάμε από αυτή τη διάλεξη}
\begin{enumerate}
  \item \textbf{Σωρός} --- ισοσταθμισμένο δυαδικό δέντρο με ιδιότητα διάταξης
  (γονιός $\le$ παιδιά). Αποθηκεύεται σε πίνακα:
  $\text{parent}(i) = \lfloor i/2 \rfloor$, παιδιά $2i$ και $2i+1$.
  \item \textbf{Heapify-up / Heapify-down} --- επιδιορθώνουν μία παραβίαση σε
  $O(\log n)$, σπρώχνοντας ένα στοιχείο προς τη ρίζα ή προς τα φύλλα (πάντα προς
  το \textbf{μικρότερο} παιδί).
  \item \textbf{Ουρά προτεραιότητας} --- \texttt{FindMin} $O(1)$,
  \texttt{Insert}/\texttt{Delete}/\texttt{ExtractMin} $O(\log n)$. Δίνει
  \textbf{Heapsort} $O(n \log n)$ και κάνει τον \textbf{Dijkstra} \& \textbf{Prim}
  $O(m \log n)$.
  \item \textbf{Union-Find} --- \texttt{MakeSet}, \texttt{Union}, \texttt{Find}
  πάνω σε διαμέριση σε ξένα σύνολα. Κάθε σύνολο $=$ κατευθυνόμενο δέντρο με ρίζα
  $=$ αντιπρόσωπο.
  \item \textbf{Union by size} κρατάει το βάθος $\le \log n$· \textbf{path
  compression} το μικραίνει κι άλλο. Μαζί $\to O(\alpha(n))$ ανά πράξη, πρακτικά
  σταθερό.
  \item Το union-find κάνει τον \textbf{Kruskal} $O(m \log n)$ --- η ταξινόμηση
  των ακμών είναι πια το ακριβό κομμάτι.
\end{enumerate}
\end{recapbox}

\lecturefooter
\end{document}
